\clearpage\begingroup
\providecommand{\ObsSec}[2]{#1\mathord{\overset{\mathrm{s}}{,}}#2}
\providecommand{\ObsArc}[2]{#1\mathord{\overset{\prime\prime}{,}}#2}
\providecommand{\ObsJupiter}{\mathord{\tikz[baseline=-.1ex,x=1em,y=1em,line width=.045em]{
  \draw (.02,.48) .. controls (.04,.70) and (.33,.69) .. (.31,.46) .. controls (.29,.29) and (.14,.19) .. (.02,.13)--(.49,.13);
  \draw (.37,.64)--(.37,-.06);}}}
\providecommand{\ObsVesta}{\mathord{\tikz[baseline=-.1ex,x=1em,y=1em,line width=.045em]{
  \draw (.04,.02) rectangle (.40,.38);
  \draw (.22,.38)--(.22,.48) (.12,.63)--(.22,.48)--(.32,.63);}}}

\hypertarget{native-280-14}{}
\HistoricalHead{275}{R}{VERSCHIEDENE BEOBACHTUNGEN IN GÖTTINGEN.}
\begin{center}[14.]\quad \textsc{Gauss} an \textsc{Encke}, Göttingen, 1. December 1816.\end{center}
\noindent\hspace*{\parindent}\mbox{. . . . . .} Auf Ihren Aufsatz\textsuperscript{[*)]} über den Cometen von 1812 freue ich mich im Voraus. Ich habe im Tagebuch [der Sternwarte] nachgesehen; viel wird es nicht sein, was Sie werden gebrauchen können. Ich habe ihn nur eine Nacht beobachtet, am 8. Sept. [1812] um $2^{\mathrm h}45^{\mathrm m}$ Sternzeit, wo ich ihn mit zwei Sternen der \textit{Hist. céleste}
\begin{center}\begin{tabular}{llrrr}
A.&p. 52&$8^{\mathrm h}$&$31^{\mathrm m}$&$\ObsSec{52}5$\\
B.&p. 216&$8^{\mathrm h}$&$33^{\mathrm m}$&$\ObsSec8{5}$
\end{tabular}\end{center}
\noindent verglichen habe:
\begin{center}\begin{tabular}{llllrr}
auf A&folgte&der&Comet&$4^{\mathrm m}$&$\ObsSec4{375}$\\
» B&»&»&»&$3^{\mathrm m}$&$\ObsSec9{75}$;
\end{tabular}\end{center}
\noindent den Declinationsunterschied habe ich noch nicht berechnet\textsuperscript{[**)]}. Wollen Sie gefälligst die scheinbaren Positionen dieser Sterne reduciren, so will ich hernach die Cometenörter mit Rücksicht auf den Werth jeder Vergleichung Ihnen genau bestimmen. Hr. Professor \textsc{Harding} hat den Cometen am 13. und 14. September mit dem kleinen Dollond beobachtet. Leider habe ich das Heft meiner Beob[achtungen], worin dessen Gesichtsfeld bestimmt ist, und welches
\begingroup\renewcommand{\thefootnote}{[*)}\footnotetext{\textsc{Lindemann} und \textsc{Bohnenberger}, Zeitschrift für Astronomie, Bd. II, S. 377 (1816).]}\endgroup
\begingroup\renewcommand{\thefootnote}{[**)}\footnotetext{Im \textit{Tagebuch der Sternwarte}, II. Bd., S. 1, ist das Beobachtungsresultat wie folgt angegeben:]
\begin{center}\begin{tabular}{llrrr}
$2^{\mathrm h}44^{\mathrm m}\ObsSec{49}9$ Sternzeit&AR.&$129^0$&$12'$&$\ObsArc{30}5$\\
&Decl.&14&54&$42{,}1$.
\end{tabular}\end{center}}\endgroup
\HistoricalFoot{}{35\textsuperscript{*}}
\clearpage\SourceLeaf{281}\HistoricalHead{276}{L}{PRAKTISCHE UND SPHÄRISCHE ASTRONOMIE. BRIEFWECHSEL.}
\noindent bei meinem Umzuge verkramt ist, noch nicht wieder auffinden können. Wollen Sie indessen vorläufig die Sterne reduciren, so will ich demnächst jene Beobachtungen auch berechnen. Leider wimmeln sie nur von Schreibfehlern, die zum Theil durch Conjecturen werden herausgebracht werden müssen.\ \mbox{. . .}

\noindent [Folgen Mittheilungen über \textsc{Hardings} Beobachtungen.]
\begin{center}[15.]\quad \textsc{Gauss} an \textsc{Encke}, Göttingen, 3. März 1817.\end{center}
\noindent\hspace*{\parindent}\mbox{. . . . . .} Die Reduction der hiesigen Cometenbeobachtungen von 1812 werden Sie bereits vor längerer Zeit von Hrn. von \textsc{Lindenau} erhalten haben\textsuperscript{[*)]}.\ \mbox{. . . . . .}
\begingroup\renewcommand{\thefootnote}{[*)}\footnotetext{Die betreffende anscheinend briefliche Mittheilung an v. \textsc{Lindenau} ist nicht vorhanden; die Beobachtungen (vgl. den vorigen Brief) sind nicht veröffentlicht.]}\endgroup
\begin{center}[16.]\quad \textsc{Gauss} an \textsc{Gerling}, Göttingen, 29. Mai 1818.\end{center}
\noindent\hspace*{\parindent}\mbox{. . . . . .} Künftig soll, hoffe ich, keine Planetenopposition von mir unbeobachtet bleiben. Für die Vesta war es diesmal zu spät. Hoffentlich haben Sie bereits andere gute Beobachtungen. Nur einmal habe ich sie beobachtet. Mein Resultat, worauf ich jedoch, da damals am Kreise noch berichtigt wurde, wenig Gewicht lege, mögen Sie einmal mit vergleichen, es ist folgendes:
\begin{center}\begin{tabular}{llrrr}
1818 April 20.\quad $11^{\mathrm h}25^{\mathrm m}\ObsSec{38}4$ M. Z.&AR. $=$&$199^0$&$47'$&$\ObsArc{28}8$\\
&D. $=$&$+5$&20&$17{,}3$.
\end{tabular}\end{center}
\noindent\hspace*{\parindent}Da der Kreis vom Meridian etwas abwich, so werden Sie die Culminationszeit von der in Sternzeit verwandelten AR. um eine Kleinigkeit verschieden finden. Gegenwärtig ist der Kreis bis auf ein Paar Zehntel einer Zeitsecunde im Meridian. In den nächsten Tagen werde ich die Beobachtungen für die Uranusopposition anfangen.
\begin{center}[17.]\quad \textsc{Gauss} an \textsc{Encke}, Göttingen, 1819 Febr. 25.\end{center}
\noindent\hspace*{\parindent}\mbox{. . . . . .} Meine Beobachtungen [der Pallas] sind eigentlich nicht zu meiner Zufriedenheit, weil ich meine zu feinen Fäden nicht gut beleuchten konnte und nur die \textls[120]{drei vollständigen} verdienen Vertrauen\textsuperscript{[**)]}
\begingroup\renewcommand{\thefootnote}{[**)}\footnotetext{Die beiden Beobachtungen vom 31. August und 4. September hat \textsc{Gauss} veröffentlicht, Astronomisches Jahrbuch für 1821, Werke VI, S. 416.]}\endgroup
\clearpage\SourceLeaf{282}\HistoricalHead{277}{R}{VERSCHIEDENE BEOBACHTUNGEN IN GÖTTINGEN.}
\begin{center}\fontsize{10}{12}\selectfont\setlength{\tabcolsep}{3pt}
\begin{tabular}{llr l rrr rrr l}
{}[1818]&Aug.&31.&$12^{\mathrm h}25^{\mathrm m}\ObsSec8{7}$&$345^0$&$49'$&$\ObsArc{45}1$&$+2^0$&$23'$&$\ObsArc2{4}$&\\
&Sept.&1.&—&\multicolumn{3}{c}{—}&$+2$&10&$22{,}9$&$\mathord{:}\,\mathord{:}\,\mathord{:}$\\
&&4.&$12\quad6\quad26{,}8$&345&5&$5{,}8$&$+1$&32&$43{,}4$&\\
&&10.&—&\multicolumn{3}{c}{—}&$+0$&14&$9{,}2$&$\mathord{:}$\\
&&12.&—&\multicolumn{3}{c}{—}&$-0$&12&$33{,}3$&$\mathord{:}\,\mathord{:}$\\
&&13.&—&\multicolumn{3}{c}{—}&$-0$&25&$43{,}0$&$\mathord{:}\,\mathord{:}$\\
&&17.&$11\quad5\quad37{,}5$&342&39&$4{,}8$&$-1$&19&$29{,}5$.&
\end{tabular}\end{center}
\begin{center}[18.]\quad \textsc{Gauss} an \textsc{Olbers}; Göttingen, 9. Oktober 1825.\end{center}
\begin{center}\small
\rule{.92\textwidth}{.35pt}\\[.2em]
\textsc{Wilhelm Olbers}, Sein Leben und seine Werke, II\textsubscript{2}, Berlin 1909, S. 428.\\[-.2em]
\rule{.92\textwidth}{.35pt}
\end{center}
\noindent\hspace*{\parindent}\mbox{. . . . . .} Ich habe ihn [den Cometen 1825 IV] zweimal im Meridian zu beobachten versucht, allein die Beobachtungen taugen nicht viel, denn so hübsch er jetzt dem blossen Auge vorkommt, so elend wird er schon bei der kleinsten 80m[aligen] Vergrösserung des Mer[idian]-Kreises und der allergeringsten Beleuchtung, die ihn ganz auslöscht.

Hier die Beobachtungen
\begin{center}\begin{tabular}{lr rrr rrr}
1825 Okt.&4.&$38^0$&$0'$&$50''$&$-14^0$&$6'$&$18''$\\
&6.&34&3&$42$\textsuperscript{[*)]}&$-18$&16&57.
\end{tabular}\end{center}
\begingroup\renewcommand{\thefootnote}{[*)}\footnotetext{Diese Zahl ist nach dem darauf folgenden Briefe von \textsc{Olbers} an \textsc{Gauss} vom 12. Oktober verbessert; vgl. auch den Brief von \textsc{Gauss} an \textsc{Schumacher} vom 21. November.]}\endgroup
\begin{center}[19.]\quad \textsc{Gauss} an \textsc{Schumacher}, Göttingen, 25. Juni 1831.\end{center}
\begin{center}\small
\rule{.92\textwidth}{.35pt}\\[.2em]
Briefwechsel zwischen \textsc{Gauss} und \textsc{Schumacher}, Altona 1860—65, II, S. 264.\\[-.2em]
\rule{.92\textwidth}{.35pt}
\end{center}
\noindent\hspace*{\parindent}\mbox{. . . . . .} Von der neulichen Sternbedeckung konnte ich wegen Wolken den Austritt nicht sehen. Den Eintritt habe ich aber sehr gut beobachtet:

Eintritt $\gamma$ Librae am dunkeln Mondsrande
\begin{center}1831 Juni 21.\quad $16^{\mathrm h}8^{\mathrm m}\ObsSec{23}4$\quad Sternzeit.\end{center}
\clearpage\SourceLeaf{283}\HistoricalHead{278}{L}{PRAKTISCHE UND SPHÄRISCHE ASTRONOMIE. BRIEFWECHSEL.}
\begin{center}[20.]\quad \textsc{Gauss} an \textsc{Schumacher}, Göttingen, 27. September 1835.\end{center}
\begin{center}\small
\rule{.92\textwidth}{.35pt}\\[.2em]
Briefwechsel zwischen \textsc{Gauss} und \textsc{Schumacher}, Altona 1860—65, II, S. 421.\\[-.2em]
\rule{.92\textwidth}{.35pt}
\end{center}
\noindent\hspace*{\parindent}\mbox{. . . . . .} Den Cometen [1835 III Halley] habe ich nach dem Mondschein zum ersten Mal wieder beobachtet:
\begin{center}\small\setlength{\tabcolsep}{2pt}
\begin{tabular}{lrrr l lrrr}
1835 Sept. 21.&$11^{\mathrm h}$&$25^{\mathrm m}$&$26^{\mathrm s}$&M. Z.&Abw[eichung] $=$&$30^0$&$33'$&$\ObsArc5{6}$\\
&11&29&25&&G[erade] A[ufsteigung] $=$&93&33&$20{,}0$
\end{tabular}\end{center}
\noindent ohne noch bei der Reduction auf die Refraction Rücksicht zu nehmen.
\clearpage\endgroup
